\documentclass[11pt]{article}
\usepackage[T1]{fontenc}
\usepackage{lmodern}
\usepackage[margin=1in]{geometry}
\usepackage{amsmath,amssymb,amsthm,booktabs,microtype}
\usepackage[hidelinks,hypertexnames=false]{hyperref}\hypersetup{pdftitle={A Newton-positivity conjecture for dyadic stages of reciprocal-subsum denominators},pdfauthor={Xiaojian Zeng},pdfsubject={Conjecture note; original log-concavity problem unresolved},pdfkeywords={subsum polynomials, log-concavity, Newton coefficients, exact certificates}}
\usepackage{url}
\newtheorem{theorem}{Theorem}
\newtheorem{proposition}[theorem]{Proposition}
\newtheorem{lemma}[theorem]{Lemma}
\newcommand{\Z}{\mathbb Z}
\newcommand{\Lfun}{\mathcal L}
\newcommand{\ind}{\mathbf 1}
\DeclareMathOperator{\odd}{o}
\title{A Newton-positivity conjecture for dyadic stages\\
of reciprocal-subsum denominators}
\author{Xiaojian Zeng\\
\small PBC School of Finance, Tsinghua University, 100190\\
\small \texttt{xiaojianzeng@foxmail.com}}
\date{}
\begin{document}
\maketitle
\begin{abstract}
The conjectured coefficient log-concavity of ordinary reciprocal-subsum
denominators remains unresolved. We propose nonnegativity of all Newton
coefficients of a quadratic minor margin at actual intermediate dyadic
stages, starting when the relevant factor exponents become affine.
An exact polynomial reduction bounds the degree by twice the coefficient
index. Complete integer certificates through index 96 imply log-concavity
at the first and last 96 internal indices for every admissible denominator.
We also prove a uniform band of nonnegative high-order differences and an
obstruction to extending the polynomial method through its first excluded
degree. Larger finite tests support the conjecture, with their verification
limits stated explicitly. Neither the proposed Newton inequalities nor
the unrestricted denominator conjecture is proved.
\end{abstract}

\section{The original problem and the proposed inequalities}

For $n\geq1$, let
\begin{equation}\label{eq:D}
 D_n(x)=\prod_{i=1}^{n}(1+x^i)^{e_{n,i}},\qquad
 e_{n,i}=\#\{q\in\Z_{\geq0}:2^qi\leq n\},\qquad L_n=\deg D_n.
\end{equation}
Ballantine, Beck, Feigon and Maurischat introduced the subsum polynomial
$\operatorname{sp}(\lambda,x)=\prod_{a\in\lambda}(1+x^a)$ of a partition
and studied sums of its reciprocals. Their Proposition~3.8 identifies
$D_n$ with their ordinary denominator $\operatorname{den}(n,x)$
\cite{BBFM}. Writing $D_n(x)=\sum_r d_{n,r}x^r$, their Conjecture~4 asks for
\begin{equation}\tag{LC}\label{eq:LC}
 d_{n,j}^2\geq d_{n,j-1}d_{n,j+1}\quad
 \left\{\begin{array}{l}
 n\geq1,\quad n\notin\{3,5,6,7\},\\
 1\leq j<L_n.
 \end{array}\right.
\end{equation}
All integer coefficient indices are included. The four exceptions are
already part of the original conjecture. Chen, Ono and Zhang retain this
ordinary-denominator question as unresolved in Table~1 of their second
version \cite{COZ}; their correction of a binary-numerator conjecture
concerns a different polynomial. We leave \eqref{eq:LC} unresolved.

The nearest structural results are palindromicity and unimodality, and
the recurrence \cite[Propositions 3.9--3.10]{BBFM}
\begin{equation}\label{eq:recurrence}
 \frac{D_n(x)}{D_{n-1}(x)}=
 \prod_{k=0}^{v_2(n)}(1+x^{2^k\odd(n)}),\qquad
 \odd(n)=n/2^{v_2(n)}\quad(n\geq2).
\end{equation}
Here $v_2(n)$ is the exponent of 2 in $n$. We process every factor in
\eqref{eq:recurrence} in increasing order of degree. Palindromicity
supplies reflection of minors, but unimodality supplies no sign for those
minors. A preservation argument must also handle factors with gaps in
their coefficient support.

For a polynomial $a(x)=\sum_r a_rx^r$, extend its coefficients by zero
outside its support and put
\[
 T_r(a)=a_r^2-a_{r-1}a_{r+1}\qquad(r\in\Z).
\]
At a stage adding $1+x^s$, write $n=s2^t$. Its old polynomial is exactly
\begin{equation}\label{eq:a}
 a_{s,t}(x)=D_{s2^t-1}(x)
       \prod_{\ell=1}^{v_2(s)}(1+x^{s/2^\ell}),
 \qquad t\geq0,\quad s2^t\geq9.
\end{equation}
The extra factors are all those preceding $1+x^s$, including when $s$ is
even. Define the stage margin
\begin{equation}\label{eq:F}
 F_{s,j}(t)=2T_j\bigl(a_{s,t}(1+x^s)\bigr)
                   -T_j(a_{s,t})-T_{j-s}(a_{s,t}).
\end{equation}
The useful stage inequality is
\begin{equation}\tag{H}\label{eq:H}
 2T_j\bigl(a_{s,t}(1+x^s)\bigr)
       \geq T_j(a_{s,t})+T_{j-s}(a_{s,t}),
\end{equation}
or equivalently $F_{s,j}(t)\geq0$.

\paragraph{Conjecture (Newton positivity at the affine start).}
For every integer $j\geq0$ and every integer $1\leq s\leq j+1$, set
\begin{equation}\label{eq:tau}
 \tau=\min\{t\in\Z_{\geq0}:s2^t\geq9,\quad 2s2^t>j+1\}.
\end{equation}
Then
\begin{equation}\tag{K}\label{eq:K}
 c_{s,j,h}:=\Delta^hF_{s,j}(\tau)\geq0
       \qquad(0\leq h\leq2j).
\end{equation}
Here $\Delta f(t)=f(t+1)-f(t)$ acts on the dyadic parameter $t$, not on
$n$. This conjecture is unproved. For $j=0$ its sole assertion is $F=1$.

The proposal concerns a simultaneous sign structure for higher differences
of actual intermediate-stage margins. This is distinct from the
individual final-polynomial minors in \eqref{eq:LC}. The results below
explain its starting point, give an infinite family of supporting signs,
and turn fixed coefficient edges into finite certification problems.
The comparison with \cite{BBFM,COZ} and bounded searches of primary
literature through 13 September 2026 found no equivalent statement;
this is not an exhaustive priority assessment. Newton interpolation itself
is standard \cite[Section 3.3(iv)]{DLMF}.

\section{Exact polynomial reduction}

\begin{proposition}\label{prop:reduction}
Fix $(s,j)$ in \eqref{eq:K}, let $m=j+1$, $A=a_{s,\tau}$, and put
$P_m(x)=\prod_{i=1}^m(1+x^i)$. For integers $u\geq0$,
\begin{equation}\label{eq:affine}
 a_{s,\tau+u}(x)\equiv A(x)P_m(x)^u\pmod{x^{m+1}}.
\end{equation}
The function $F_{s,j}(\tau+u)$ agrees with a polynomial in $u$ of degree
at most $2j$. In particular,
\begin{equation}\label{eq:newton}
 F_{s,j}(\tau+u)=\sum_{h=0}^{2j}c_{s,j,h}\binom uh.
\end{equation}
For $\Lfun_j(W)=[x^jy^j]W-[x^{j-1}y^{j+1}]W$, these coefficients are
\begin{equation}\label{eq:bivariate}
 c_{s,j,h}=\Lfun_j\!\left(
 A(x)A(y)\bigl(1+2x^s+2y^s+(xy)^s\bigr)
       \bigl(P_m(x)P_m(y)-1\bigr)^h\right).
\end{equation}
Coefficient extraction at a negative exponent is understood to be zero.
\end{proposition}

\begin{proof}
Let $E_i(n)=\#\{q\geq0:2^qi<n\}$, the exponent of $1+x^i$ in
$D_{n-1}$. Splitting off $q=0$ gives
\[
 E_i(2n)-E_i(n)=\ind_{i<2n}.
\]
The choice of $\tau$ makes this difference equal to 1 for every $i\leq m$
at every subsequent doubling. The extra factors in \eqref{eq:a} do not
depend on $t$. This proves \eqref{eq:affine}.

The coefficient of $x^r$ in $AP_m^u$ is a polynomial in $u$ of degree at
most $r$, by the binomial expansion of the factors. On expanding
\eqref{eq:F}, each nonzero product of old coefficients has total index
at most $2j$. Hence the margin has degree at most $2j$. The polynomial
Newton formula now gives \eqref{eq:newton}. For completeness, the basis
$\binom uh$ satisfies $\Delta\binom uh=\binom u{h-1}$, so successive
differences at zero recover its coefficients. Thus no fitted degree
assumption is involved.

Finally $T_j(a)=\Lfun_j(a(x)a(y))$, and twice the new product minus the
old product and its shift has multiplier
\[
 2(1+x^s)(1+y^s)-1-(xy)^s=1+2x^s+2y^s+(xy)^s.
\]
Each difference multiplies the remaining expression by
$P_m(x)P_m(y)-1$, proving \eqref{eq:bivariate}.
\end{proof}

If \eqref{eq:K} holds, \eqref{eq:newton} gives \eqref{eq:H} for affected
shifts $s\leq j+1$ throughout the affine region $2n>j+1$, because
$\binom uh\geq0$ for integers $u\geq0$. It supplies no automatic sign
before $\tau$. For unaffected shifts $s\geq j+2$, the coefficient triple
at $j$ is unchanged and $F_{s,j}=T_j(a_{s,t})$; this requires an old-minor
bound. Thus \eqref{eq:K} alone is not a complete induction on a moving
linear edge and is not asserted to be equivalent to \eqref{eq:LC}.

\section{A complete certificate for a fixed edge}

\begin{theorem}[Computer-assisted fixed-edge result]\label{thm:edge}
For every $n\geq9$, every actual stage in \eqref{eq:recurrence} satisfies
\eqref{eq:H} at $0\leq j\leq96$ and at their reflections in the degree
of the new stage. For every $n\geq1$ outside $\{3,5,6,7\}$,
\[
 T_j(D_n)\geq0\quad\text{if}\quad
 1\leq j<L_n,\qquad \min(j,L_n-j)\leq96.
\]
\end{theorem}

\begin{proof}
The exact finite certificate consists of all pairs
$0\leq j\leq96$, $1\leq s\leq j+1$, at their least starts
\eqref{eq:tau}, and all orders $0\leq h\leq2j$. These are 4,753 pairs
and 613,137 strictly positive integer coefficients. It also verifies all
3,126 triples $(s,j,t)$ in the same pair domain with $s2^t\geq9$ and
$t<\tau$. Every one of their margins is nonnegative. The construction
and separate recomputation are specified below.

For affected shifts, Proposition~\ref{prop:reduction} and the earlier
triples therefore prove \eqref{eq:H} for every admissible $t$.
Rebuilding the full products verifies log-concavity for $D_1,D_2,D_4,D_8$.
Starting from $D_8$, induct over every ascending stage. The old indices
$j$ and $j-s$ are at most 96; a negative index has zero minor. In an
affected stage, their nonnegativity and \eqref{eq:H} give the new minor's
nonnegativity. In an unaffected stage, the minor is unchanged, and its
old sign also proves \eqref{eq:H}. This simultaneously supplies the
old-minor hypothesis, without circularity.

All stage polynomials are palindromic. If the old degree is $L$, then
$T_r(a)=T_{L-r}(a)$, including zero-extended indices, and
$F_{s,j}=F_{s,L+s-j}$: reflection interchanges the two old minors.
This proves the reflected stage statement and the right-edge result.
\end{proof}

The finite part can be checked using only integer arithmetic. At each
fixed $s$, construct coefficients through degree 97 at each needed actual
$t$, form \eqref{eq:F}, and take forward differences. An independent
construction uses logarithmic derivatives. More explicitly, if
$a(x)=\prod_i(1+x^i)^{f_i}$, then
\begin{equation}\label{eq:logderiv}
 \beta_r=\sum_{i\mid r} i f_i(-1)^{r/i-1},\qquad
 a_0=1,\qquad r a_r=\sum_{k=1}^r\beta_k a_{r-k}.
\end{equation}
Here $f_i=E_i(s2^t)+\ind_{i\in\{s/2,\ldots,s/2^{v_2(s)}\}}$;
every division by $r$ is checked exact. Formula \eqref{eq:logderiv}
follows by differentiating the product, so it does not assume the affine
reduction. Each difference can also be reconstructed as
\begin{equation}\label{eq:signed}
 c_{s,j,h}=\sum_{u=0}^h(-1)^{h-u}\binom hu F_{s,j}(\tau+u).
\end{equation}

The original producer uses sparse shift-and-add multiplication and the
exact doubling identity. Separate programs reconstruct all 613,137
coefficients by \eqref{eq:logderiv} and \eqref{eq:signed}, as well as all
earlier triples. A further verification uses the telescoped Euler product
described in Appendix~\ref{app:reproduction}. The edge theorem uses cutoff
96 throughout; later tests below are evidence about the conjecture.

The use of a nonnegative polynomial basis has a related verification
precedent: Zeng, Liu, Ratnavelu, Huat and Xiong report exact Bernstein
certificates with separate collocation reconstruction for a Hermitian
scalar inequality \cite{ZengGate}. Here the basis is instead
$\{\binom uh\}$ on the nonnegative integers, and the infinite range follows
from the proved degree bound and the exact finite coefficients.

\section{A uniform band of high-order coefficients}

The following elementary smoothing argument applies more generally than
the particular stage polynomials.

\begin{lemma}\label{lem:smoothing}
Let $V(x,y)$ be symmetric and homogeneous of degree $d$, with nonnegative
real coefficients. If $N\geq0$, $N+d=2j$, and $N+2\geq d^2$, then
$\Lfun_j((x+y)^NV)\geq0$.
\end{lemma}

\begin{proof}
Decompose $V$ into nonnegative multiples of
$x^ay^{d-a}+x^{d-a}y^a$, allowing half a pair on the diagonal.
Set $q=j-a$. The assumptions imply $N\geq d$: this is immediate for
$d=0$, follows from parity for $d=1$, and follows from $N\geq d^2-2$
for $d\geq2$. Thus $0\leq q\leq N$. On this monomial pair the functional
is
\[
 2\binom Nq-\binom N{q-1}-\binom N{q+1}
 =\binom Nq\frac{N+2-(N-2q)^2}{(q+1)(N-q+1)}.
\]
Since $|N-2q|=|2a-d|\leq d$, the numerator is nonnegative and the
denominator positive. Binomial coefficients outside $[0,N]$ are zero;
the formula also covers $N=d=0$.
\end{proof}

\begin{proposition}\label{prop:band}
Let $B$ and $R$ be symmetric formal power series with nonnegative real
coefficients. Suppose the constant term of $R$ is zero and its linear
part is $x+y$. If integers $j,k\geq0$ satisfy
\[
 h=2j-k\geq0,\qquad j\geq2k^2+k-1,
\]
then $\Lfun_j(BR^h)\geq0$. In particular, $c_{s,j,2j-k}\geq0$ under
these conditions for every pair in \eqref{eq:K}.
\end{proposition}

\begin{proof}
Write $R=(x+y)+S$, where $S$ has no terms of degree below 2. Expand
\[
 BR^h=\sum_{r=0}^h\binom hr(x+y)^{h-r}BS^r.
\]
In the $r$th term, only the homogeneous component of $BS^r$ of degree
$d=2j-(h-r)=k+r$ contributes. It vanishes unless $d\geq2r$, so
$r\leq k$ and $d\leq2k$. Put $N=h-r=2j-d$. Then
\[
 N+2-d^2=2j+2-d(d+1)
       \geq2j+2-2k(2k+1)\geq0.
\]
Apply Lemma~\ref{lem:smoothing} to each contributing component. Only
finitely many components occur. For the application take
$B=A(x)A(y)(1+2x^s+2y^s+(xy)^s)$ and
$R=P_m(x)P_m(y)-1$ in \eqref{eq:bivariate}.
\end{proof}

There are also explicit formulas for the two highest orders. With
$C_j=\binom{2j}j/(j+1)$ and $a_1=[x]A$, one has
\begin{align}
 c_{s,j,2j}&=C_j,\label{eq:top}\\
 c_{s,j,2j-1}&=(a_1+2\ind_{s=1})C_j
              +(2j-1)(C_j-C_{j-1})\qquad(j\geq1).\label{eq:nexttop}
\end{align}
Indeed $B_0=1$, $B_1=(a_1+2\ind_{s=1})(x+y)$, and, for $j\geq1$,
$R_2=(x+y)^2-xy$. For order $2j$ only $(x+y)^{2j}$ contributes.
For order $2j-1$ the two contributions are
$B_1(x+y)^{2j-1}$ and $(2j-1)(x+y)^{2j-2}R_2$.
Using $\Lfun_j((x+y)^{2j})=C_j$ and
$\Lfun_j(xy(x+y)^{2j-2})=C_{j-1}$ proves the formulas.

The low and middle orders do not follow from coefficient positivity
alone. For example, take $A=1$, $s=3$, $j=2$ and $m=3$ in
\eqref{eq:bivariate}. Direct expansion gives $\Lfun_2(BR)=-3$ and
$\Lfun_2(BR^2)=-1$. This $A$ is not an actual old stage. The example
invalidates the corresponding generic positivity argument, while
leaving \eqref{eq:K} untouched. It explains why the actual exponent
profile in \eqref{eq:a} is a substantive hypothesis.

\section{The first excluded degree}

The strict inequality in \eqref{eq:tau} has an algebraic reason.
At its boundary, the actual sequence has a discrepancy at its first
term relative to its polynomial tail.

\begin{proposition}\label{prop:boundary}
Let $n=s2^t\geq9$ be any actual stage, put $j=2n-1$, and write
$a=a_{s,t}$. There is a polynomial $Q(u)$ of degree at most $4n-2$ such
that $F_{s,j}(t+u)=Q(u)$ for every integer $u\geq1$, whereas
\begin{equation}\label{eq:discrepancy}
 F_{s,j}(t)-Q(0)=-a_{2n-2}-2a_{2n-s-2}<0.
\end{equation}
For every integer $h>4n-2$,
\begin{equation}\label{eq:boundarydiff}
 \Delta^hF_{s,j}(t)=(-1)^{h+1}(a_{2n-2}+2a_{2n-s-2}).
\end{equation}
In particular all even differences of order at least $4n$ are negative.
\end{proposition}

\begin{proof}
Set $m=2n=j+1$. The exponent identity in the proof of
Proposition~\ref{prop:reduction} gives, for integers $u\geq1$,
\[
 a_{s,t+u}(x)\equiv\frac{a(x)P_m(x)^u}{1+x^m}
                         \pmod{x^{m+1}}.
\]
The first doubling adds every factor through degree $m-1$, while later
doublings add every factor through $m$. Coefficients of degree $r$ in
this expression are polynomial in $u$ of degree at most $r$, including
the exponent $u-1$ at degree $m$. Substitution into \eqref{eq:F} defines
$Q$ with degree at most $2j=4n-2$.

At $u=0$ the extrapolated old series is
\[
 \widehat a(x)=\frac{a(x)}{1+x^{2n}}
              \equiv a(x)-x^{2n}\pmod{x^{2n+1}}.
\]
Only $a_{j+1}$ changes among the coefficients needed for the margin.
Expanding that margin gives
\begin{align*}
 F={}&T_j(a)+T_{j-s}(a)+4a_ja_{j-s}\\
     &\ -2a_{j-1}a_{j-s+1}-2a_{j+1}a_{j-s-1}.
\end{align*}
The coefficient of $a_{j+1}$ is $-a_{j-1}-2a_{j-s-1}$, proving the
equality in \eqref{eq:discrepancy}. The factor
$\prod_{i=1}^{n-1}(1+x^i)$ divides $D_{n-1}$ and has positive
coefficients throughout $[0,n(n-1)/2]$: the integers $1,\ldots,n-1$
have every subset sum in this interval. Since $2n-2\leq n(n-1)/2$
for $n\geq9$, we have $a_{2n-2}>0$.

The actual sequence differs from $Q$ only at $u=0$. A difference of
order $h>4n-2$ annihilates $Q$; the remaining term is $(-1)^h$ times
the discrepancy. This proves \eqref{eq:boundarydiff}.
\end{proof}

This result obstructs extending the polynomial reduction and
nonnegativity of all differences unchanged to $2n=j+1$. It does not
assert a negative coefficient of order at most $2j$ at that excluded
start, nor does it assert that the actual margin is negative. For
example, at $n=s=9$ and $j=17$, exact calculation gives
$F=19668$, $Q(0)=20795$, and $\Delta^{36}F=-1127$.
Thus the obstruction is neither a counterexample to \eqref{eq:K} nor a
disproof of \eqref{eq:H} or \eqref{eq:LC}.

\section{Evidence, limitations, and a next falsification test}

Table~\ref{tab:tests} records the exact Newton-coefficient tests. A
template is a pair $(s,j)$, always with its prescribed least start and
every order $0\leq h\leq2j$. The conjecture was fixed after the design
range; the next contiguous block tested it without changing its statement.
All coefficients in the three contiguous blocks were strictly positive.
The sparse block also had no negative coefficients.

\begin{table}[ht]
\centering
\small
\begin{tabular}{lrrl}
\toprule
Domain in $j$ & Templates & Coefficients & Separate recomputation\\
\midrule
$0$--$64$ (design) & 2,145 & 185,185 & All coefficients\\
$65$--$96$ & 2,608 & 427,952 & All coefficients\\
Listed sparse degrees & 128 & 69,632 & Partial; see text\\
$97$--$128$ (later test) & 3,632 & 826,288 & Separate Euler product\\
\bottomrule
\end{tabular}
\caption{Exact finite evidence for \eqref{eq:K}. The last two rows
overlap and must not be added as disjoint tests.}\label{tab:tests}
\end{table}

The sparse degrees are
\[
 \{126,127,128,129,190,191,192,193,
   254,255,256,257,510,511,512,513\},
\]
with shifts
\[
 \{1,2,3,\lfloor(j+1)/2\rfloor,\lfloor(j+1)/2\rfloor+1,j-1,j,j+1\}
\]
at each degree. They target changes in the dyadic start, small shifts,
threshold shifts, and the onset of overlap. Their producer uses exact
integer polynomials in python-flint 0.9.0 \cite{PythonFlint}; all forward
differences use Python integers. Separate arithmetic checks the first
five orders, both highest-order formulas, and two evaluations outside
the interpolation grid for each of these 128 pairs. It does not
recompute every middle coefficient of this sparse block. No exhaustive
coverage through 513 is asserted.

The later separate Euler-product test covers every shift at
$97\leq j\leq128$. It overlaps the sparse block in 24 templates and
6,120 coefficients; hence 3,608 templates and 820,168 coefficients were
previously unlisted. These are further falsification data, without an
increase to the fixed-edge theorem's stated cutoff. The boundary
identity was initially checked at all 48 actual stages for $9\leq n\leq32$,
and subsequently at all 112 stages for $9\leq n\leq64$.

The unbounded low and middle Newton orders remain uncontrolled.
Theorem~\ref{thm:edge} concerns a fixed distance from either edge; it
leaves the unbounded interior of \eqref{eq:LC} unresolved. The arbitrary
cutoff 96 is a certified range, not an extremal threshold. The
conjectural sign structure, the high-order band, and the precise
boundary obstruction together are the proposed contribution.

A reproducible falsifier for \eqref{eq:K} is any admissible tuple
$(s,j,h)$ with the unique $\tau$ of \eqref{eq:tau} and a negative exact
sum \eqref{eq:signed}. One next test is
$j\in\{1022,1023,1024\}$, all
$s\in\{1,2,3,511,512,513,j,j+1\}$, and all $0\leq h\leq2j$.
A negative coefficient must be distinguished from a negative actual
margin: the former refutes \eqref{eq:K} but need not refute
\eqref{eq:H}. Only a negative final-polynomial minor at an admissible
$n$ would refute \eqref{eq:LC}.

\appendix
\section{Reproduction of the fixed certificate}\label{app:reproduction}

The accompanying \texttt{verify\_edge.py} reproduces exactly the
$j\leq96$ certificate and the finite bases, reading the supplied data
without changing them. It is a presentation adaptation of the separate
Euler-product verifier, not an additional independent implementation.
It requires only Python~3 standard-library integer arithmetic.

For $N=s2^t-1$ and $o=\odd(s)$, dyadic telescoping gives
\begin{equation}\label{eq:euler}
 a_{s,t}(x)=\frac{1-x^s}{1-x^o}
 \prod_{\substack{d\leq N\\d\ \mathrm{odd}}}
       (1-x^d)^{-e_{N,d}}
 \prod_{\substack{d\leq2N\\d\ \mathrm{even}}}(1-x^d).
\end{equation}
To see this, replace each $1+x^i$ by $(1-x^{2i})/(1-x^i)$ in $D_N$.
At an even index $d\leq N$, the difference between numerator and
denominator exponents is $e_{N,d/2}-e_{N,d}=1$; for
$N<d\leq2N$ the numerator exponent is also 1. At an odd index only
the denominator remains. The extra stage factors telescope to
$(1-x^s)/(1-x^o)$.

Writing \eqref{eq:euler} as $\prod_d(1-x^d)^{-g_d}$ gives
$\beta_r=\sum_{d\mid r}d g_d$. The same coefficient recursion in
\eqref{eq:logderiv} then reconstructs the old polynomial at every
actual $t$. For each $s=1,\ldots,97$, take all $t$ from its first value
with $s2^t\geq9$ through $\tau(s,96)+192$. Compute the margin vector
for $j=0,\ldots,96$ at these $t$, select the $2j+1$ values beginning at
$\tau(s,j)$, and form their difference table. Compare every entry with
the saved coefficients, check its sign, and check the exact earlier
triple domain. Full direct products verify the four finite bases.
The reproduction README identifies the data files and commands.

\paragraph{Computational and AI assistance.}
OpenAI Codex assisted with conjecture formulation, algebraic arguments,
programming, literature checks, and manuscript preparation. A separate
model context reconstructed the arguments and implemented the
Euler-product verification. Separate contexts and separate programs do
not constitute statistically independent evidence or human peer review.
The supporting computer-assisted theorem uses the stated exact integer
checks and polynomial identities; no floating-point sign decisions,
random sampling, or proof-assistant formalization are used. The sparse
tests use python-flint only as described above.


\small
\raggedright
\bibliographystyle{plain}
\bibliography{references}

\end{document}
